\documentclass{beamer}
\usetheme{Pittsburgh}

\usepackage{enumitem, colortbl}
\usepackage[pagewise,displaymath, mathlines]{lineno}
\usepackage{comment}
\usepackage{listings}
\usepackage{fancyvrb}


%\usepackage{tikz}

%\usetikzlibrary{automata,positioning}

\usepackage{animate}
%\usepackage[scaled]{berasans}
\renewcommand*\familydefault{\sfdefault} %% Only if the base font of the document is to be sans serif

\usepackage[T1]{fontenc}

%\usepackage{epstopdf}
\setbeamertemplate{navigation symbols}{}
\setbeamertemplate{background}{\includegraphics[width=\paperwidth]{../Brown-Paper-6.jpg}}

\title % (optional, only for long titles)
{\textbf{Caret and zoon}: machine learning, ecology and domain specific package systems}
\author % (optional, for multiple authors)
{Tim~\small{C.D.}~\normalsize{Lucas}}
\institute % (optional)
{
  Malaria Atlas Project\\
  Oxford
}
%\usepackage{PSTricks}

\makeatletter
\newlength\beamerleftmargin
\setlength\beamerleftmargin{\Gm@lmargin}


\begin{document}

\begin{frame}

\frametitle{\textbf{Caret and zoon}: machine learning, ecology and domain specific package systems
\\}
\vspace{3cm}
\framesubtitle{{\color{black} \normalsize{Tim}~\small{C.D.}~\normalsize{Lucas}
\hbox{}\\
\vspace{3mm}
{\small Malaria Atlas Project, BDI, Oxford}}\\
\includegraphics[height=7pt]{Twitter_logo_blue-small.png}{\color{black} \hspace{0.6mm}\footnotesize{@}\small{timcdlucas}  \footnotesize{@}\small{statsforbios}}\\
\includegraphics[height=7pt]{Ar_Icon_Contact.pdf}{\color{black}  \hspace{0.6mm}\small{timcdlucas}\footnotesize{@}\small{gmail.com}}}

%\hspace{-6cm}
\vspace{-3cm}
%\hbox{
%\include{gasModel1.tex}
%\includegraphics[height=0.5\textheight]{gasModel1.png}
%}
\end{frame}




\begin{frame}
\frametitle{Who am I?}
	\begin{itemize}
	\item Malaria Atlas Project at BDI
	\item Malaria, maps, geostatistics

	\end{itemize}	
% Your image included here
	\centerline{
    \includegraphics[width=0.8\paperwidth]{03_Country_Comparison_API_India.png}
	}
\end{frame} 




\begin{frame}
\frametitle{Who am I?}
\framesubtitle{R packages}
	\begin{itemize}
	\item Zoon 
	\item INLAutils
	\item palettetown - my greatest ever achievement
	\end{itemize}	
% Your image included here
	\centerline{
    \includegraphics[height = 0.59\textheight]{pokedex-4.png}
	}
\end{frame} 




\begin{frame}
\frametitle{Talk overview}
\begin{block}{caret}
	\begin{itemize}
	\item General package for machine learning.
	\item Introduction to the package.
    \item A domain specific package ecosystem?
	\end{itemize}
\end{block}
\begin{block}{zoon}
	\begin{itemize}
		\item General package for species distribution modelling.
    \item What are SDMs?
    \item Package overview.
	\end{itemize}

\end{block}



\begin{block}{Domain specific ecosystems}

	\begin{itemize}
	\item Other examples.
	\item Are they a good thing?
	\end{itemize}
\end{block}

\end{frame} 












\begin{frame}
\frametitle{What is machine learning?}
\begin{figure}
    \includegraphics[width = 0.6\textwidth]{classification}
\end{figure} 
\begin{figure}

    \includegraphics[width = 0.6\textwidth]{regression}
\end{figure} 
\end{frame} 




\begin{frame}
\frametitle{What is machine learning?}
\begin{itemize}
\item Only care about prediction
\item Not mechanistic/process based models
\item Not inference
\end{itemize}
\end{frame} 



\begin{frame}
\frametitle{Cross-validation}
\begin{figure}
 \animategraphics[loop,autoplay, width = 0.6\textwidth]{1}{frame-}{0}{2}
\end{figure} 
\end{frame} 


\begin{frame}
\frametitle{Hyperparameters}
\begin{figure}
 \animategraphics[loop,autoplay, width = 0.7\textwidth]{1}{scale-}{0}{2}
\end{figure} 
\end{frame} 


\begin{frame}
\frametitle{Hyperparameters}
\begin{itemize}
\item Number of PCA coordinates
\item Cut-offs for variable selection
\item $x + x^2 + x^3 + x^4 + ...$
\end{itemize}
\end{frame} 



\begin{frame}
\frametitle{No free lunch}
No such thing as a universal, `best' machine learning model.
\end{frame} 



\begin{frame}
\frametitle{caret}
https://topepo.github.io/caret/model-training-and-tuning.html
\end{frame} 



\begin{frame}
\frametitle{What does caret do?}
\begin{figure}
    \includegraphics[height = 0.8\textheight]{Capture}
\end{figure} 
\end{frame} 



\begin{frame}
\frametitle{What does caret do?}
\begin{figure}
    \includegraphics[height = 0.8\textheight]{gbm2opt}
\end{figure} 
\end{frame} 






\begin{frame}[fragile]
\frametitle{Training a model}
\begin{Verbatim}

m1 <- train(Species ~ ., 
            iris,
            method = `gbm')

\end{Verbatim}

\begin{figure}
    \includegraphics[height = 0.5\textheight]{train_gbm}
\end{figure} 
\end{frame} 


\begin{frame}[fragile]
\frametitle{Training a different model}
\renewcommand{\FancyVerbFormatLine}[1]{%
   \ifnum\value{FancyVerbLine}=3\color{cyan}#1%
   \else #1\fi}

\begin{Verbatim}
m2 <- train(Species ~ ., 
            iris,
            method = `nnet')
\end{Verbatim}

\begin{figure}
    \includegraphics[height = 0.5\textheight]{train_nnet}
\end{figure} 
\end{frame} 




\begin{frame}[fragile]
\frametitle{Controlling Crossvalidation}
\renewcommand{\FancyVerbFormatLine}[1]{%
   \ifnum\value{FancyVerbLine}=1\color{cyan}#1%
   \else #1\fi}

\begin{Verbatim}
tr <- trainControl(method = `cv', number = 5)

m3 <- train(Species ~ ., 
            iris,
            trControl = tr,
            method = `nnet')
\end{Verbatim}

\end{frame} 



\begin{frame}[fragile]
\frametitle{Try more hyperparameter values}
\renewcommand{\FancyVerbFormatLine}[1]{%
   \ifnum\value{FancyVerbLine}=3\color{cyan}#1%
   \else #1\fi}

\begin{Verbatim}
m4 <- train(Species ~ ., 
            iris,
            tuneLength = 10,
            method = `nnet')
\end{Verbatim}

\begin{figure}
    \includegraphics[height = 0.5\textheight]{train_nnet_tuneLength}
\end{figure} 
\end{frame} 





\begin{frame}[fragile]
\frametitle{Use chosen hyperparameter values}
\renewcommand{\FancyVerbFormatLine}[1]{%
   \ifnum\value{FancyVerbLine}=3\color{cyan}#1%
   \else%
   \ifnum\value{FancyVerbLine}=4\color{cyan}#1%
   \else #1\fi\fi}

\begin{Verbatim}
m5 <- train(Species ~ ., 
            iris,
            tuneGrid = expand.grid(size=c(1,5,10,20), 
                                   decay=seq(0,0.1,0.01)),
            method = `nnet')
\end{Verbatim}

\begin{figure}
    \includegraphics[height = 0.5\textheight]{train_nnet_tuneGrid}
\end{figure} 
\end{frame} 





\begin{frame}
\frametitle{Contributions}
	\begin{itemize}
	\item Add your own models.
	\item Share by github pull request.
	\item But aim is for devs to keep package up to date.
	\end{itemize}
\end{frame} 




%%%%%%%%%%%%%%%%%%%%%%%%%%%%%%%%%%%%%%%%%%%%%%%%%%%%%%%%%%%%%%%%%%%%%%%%%%%%%%%%


%%%% zoon.




\begin{frame}[plain]
    \hspace*{-\beamerleftmargin}%
    \includegraphics[width=\paperwidth,width=\paperwidth]{zoonSlide}
\end{frame} 



\begin{frame}
\frametitle{Who develops zoon?}
\begin{itemize}
\item Tom August
\item Me
\item Nick Golding
\item Emiel van Loon
\item David Gavaghan
\item Greg McInerny
\end{itemize}
\end{frame} 

\begin{frame}
\frametitle{What is species distribution modelling?}
\begin{figure}
    \includegraphics[height = 0.6\textheight]{classification2}
\end{figure} 
\end{frame} 


\begin{frame}
\frametitle{What is species distribution modelling?}
\begin{figure}
    \includegraphics[height = 0.8\textheight]{fengChangeWorkflow-1.pdf}
\end{figure} 
\end{frame} 


\begin{frame}[fragile]
\frametitle{A basic workflow}
%\vspace{-1cm}
\begin{figure}
    \includegraphics[width = \textwidth]{diagrams}
\end{figure} 
\begin{Verbatim}
work1 <- workflow(
  occurrence = UKAnophelesPlumbeus,
  covariate  = UKAir,
  process    = OneHundredBackground,
  model      = RandomForest,
  output     = PrintMap
)

\end{Verbatim}

\end{frame} 


\begin{frame}[fragile]
\frametitle{What does zoon do?}

\begin{figure}
    \includegraphics[height = 0.85\textheight]{modules}
\end{figure} 
\end{frame}




\begin{frame}[fragile]
\frametitle{A different model}
%\vspace{-1cm}
\begin{figure}
    \includegraphics[width = \textwidth]{diagrams}
\end{figure} 
\renewcommand{\FancyVerbFormatLine}[1]{%
   \ifnum\value{FancyVerbLine}=5\color{cyan}#1%
   \else #1\fi}

\begin{Verbatim}
work2 <- workflow(
  occurrence = UKAnophelesPlumbeus,
  covariate  = UKAir,
  process    = OneHundredBackground,
  model      = MaxEnt,
  output     = PrintMap
)

\end{Verbatim}

\end{frame} 



\begin{frame}[fragile]
\frametitle{A different workflow}
%\vspace{-1cm}
\begin{figure}
    \includegraphics[width = \textwidth]{diagrams}
\end{figure} 
\renewcommand{\FancyVerbFormatLine}[1]{%
   \ifnum\value{FancyVerbLine}=3\color{cyan}#1%
   \else%
   \ifnum\value{FancyVerbLine}=6\color{cyan}#1%
   \else #1\fi\fi}

\begin{Verbatim}
work3 <- workflow(
  occurrence = UKAnophelesPlumbeus,
  covariate  = UKBioclim,
  process    = Background(n = 500),
  model      = RandomForest,
  output     = Appify
)

\end{Verbatim}

\end{frame} 


\begin{frame}[fragile]
\frametitle{caret in zoon}
\begin{figure}
    \includegraphics[width = \textwidth]{diagrams}
\end{figure} 
\renewcommand{\FancyVerbFormatLine}[1]{%
   \ifnum\value{FancyVerbLine}=5\color{cyan}#1%
   \else%
   \ifnum\value{FancyVerbLine}=6\color{cyan}#1%
   \else #1\fi\fi}
\begin{Verbatim}
work4 <- workflow(
  occurrence = UKAnophelesPlumbeus,
  covariate  = UKAir,
  process    = OneHundredBackground,
  model      = MachineLearn(method = `nnet', 
                            tuneLength = 8),
  output     = PrintMap
)
\end{Verbatim}

\end{frame} 




\begin{frame}[fragile]
\frametitle{A more complicated workflow}
%\vspace{-1cm}
\begin{figure}
    \includegraphics[width = \textwidth]{diagrams}
\end{figure} 
\begin{Verbatim}
work5 <- workflow(
  occurrence = SpOcc(species = `Eresus kollari', 
                     extent = c(-10, 10, 45, 65)),
  covariate  = UKBioclim,
  process    = BackgroundAndCrossvalid(k = 5),
  model      = list(LogisticRegression,
                    RandomForest),
  output     = Chain(PrintMap(plot = FALSE),
                     PerformanceMeasures)
)
\end{Verbatim}

\end{frame} 




\begin{frame}
\frametitle{Contributions}
	\begin{itemize}
	\item Add your own methods.
	\item Share by web form or github.
	\item Not the aim for devs to keep package up to date.
	\end{itemize}
\end{frame} 




%%%%%%%%%%%%%%%%%%%%%%%%%%%%%%%%%%%%%%%%%%%%%%%%%%%%%%%%%%%%%%%%%%%%%%%%%%%%%%%%


%%%% package ecosystems.



\begin{frame}
\frametitle{Package ecosystems}
\huge{Cue hand-waving}
\end{frame} 

\begin{frame}
\frametitle{Package ecosystems}
\begin{figure}
    \includegraphics[height = 0.8\textheight]{continuum.pdf}
\end{figure} 
\end{frame} 



\begin{frame}
\frametitle{Package ecosystems}
\begin{figure}
    \includegraphics[height = 0.8\textheight]{continuum_traits.pdf}
\end{figure} 
\end{frame} 




\begin{frame}
\frametitle{Any Questions ?}
\vspace{2cm}
\framesubtitle{{\color{black} \normalsize{Tim}~\small{C.D.}~\normalsize{Lucas}}}
\end{frame}
% title
% Talk structure
% Why
% gas model pic1
% gas model pic2
% gas model pic3
% gas model equation?
%  
%
%
% 
% Model error
% Model error square
% 
%
%
%
%
%
%


\end{document}
